\scanpage{9}

\unsure{Lodève} \unsure{le}\unsure{.|,|2} 7.8.74\\
Cher Deligne,\\

Étant peut-être empêché par \unsure{ma|une} jambe d'assurer un cours de $1^{\mathrm{er}}$ cycl\unsure{\phantom{e}|e} au $1^{\mathrm{er}}$ trimestre, je vais peut-être \unsure{a|à|c} la place faire un petit séminaire d'algèbre, et envisage de le faire sur les fourbis de \unsure{Mme} Sinh, \unsure{éventuellement} transposés dans le contexte des ``champs''. À ce propos, je tombe sur le truc suivant,
% margin note written sideways, bottom-to-top; no attachment sign is visible
\margin{\unsure{je n'ai pas} \unsure{retrouvé} la \unsure{lettre} \unsure{sur les} \unsure{Gr-Champs|Gr-champs} que tu \unsure{cites|cités}}
qui pour l'instant reste heuristique. Si $M,N$ sont deux faisceaux abéliens sur un topos $X$, et $\tau_{\leq 2}\mathbb{R}\underline{\mathrm{Hom}}(M,N)=E(M,N)$ est le complexe ayant les invariants
\[
\left\{\begin{aligned}
\underline{H}^{i}&=\underline{\deleted{\unsure{\mathrm{Hom}|\mathrm{Ext}}}\mathrm{Ext}}^{i}(M,N) &&\text{pour }0=i\leq 2\\
\underline{H}^{j}&=0 &&\text{si }i\notin[0,2]
\end{aligned}\right.
\]
il doit y avoir un triangle distingué canonique
\diagram{\begin{tikzcd}[ampersand replacement=\&,column sep=small]
{} \& {} \& \underline{\unsure{\mathrm{Hom}|\mathrm{Hm}}}(M,{}_{2}N)[-2]\arrow[dl] \& {}\\
(T) \& E(M,N)\arrow[rr] \& {} \& E'(M,N)\arrow[ul]
\end{tikzcd}}

% archive sheet number: 8
